\begin{proof}

\begin{enumerate}[label={$\langle1\rangle$\arabic*.},leftmargin=*]
\item We check the root-curve and logarithmic tangent inputs of \SeriesRef{fam:src:A.1.1} and \SeriesRef{fam:st:A.5.2}. The root charts at $p_1,p_2$ are the finite etale torus-family covers used in the construction; finite-group invariants over $\C$ are exact. At the cusp the local equations are $xy=q$ and $xyz=q$, and the vertical logarithmic fields are generated by $x\partial_x-y\partial_y$ and the remaining relative direction. Their affine transition derivatives depend only on the base. The invariant vertical field gives a trivial subline of the vertical bundle $E$. The quotient has twelfth power $\OO(-p_0)$: its automorphy is that of the reciprocal elliptic discriminant, whose only divisor is a simple pole at the cusp. As $\Pic(\mathfrak B)=\Z$, the quotient is $\OO(-1)$; $H^1(\mathfrak B,\OO(1))=0$ splits the extension. Thus $E=\OO\oplus\OO(-1)$.

\item The normalization complex of $W$ has global scalar complex
\[
 \C\xrightarrow{0}\C^3
 \xrightarrow{(a,b,c)\mapsto(a-b+c,a-b+c)}\C^2.
\]
The rational normalization and the three rational double curves have no positive-degree structure-sheaf cohomology. The displayed complex therefore gives $(1,2,1)$ for the fibre cohomology and its exterior cup product; the two oriented triangular faces give the nondegenerate degree-one pairing. The cyclic-cover fibres have the same dimensions. Proper flat base change on these charts gives locally free $G=R^1\bar f_*\OO_X$, with $R^2\bar f_*\OO_X=\det G$. Relative duality and the determinant of the logarithmic sequence give $\det G=\OO(-1)$. The integral class $\gamma$ maps to a nowhere-zero section of $G$: on a marked smooth torus the two identity columns of the period matrix prevent it from belonging to the span of the two holomorphic period rows, and at the cusp it is a nonzero cycle in the displayed complex. The same splitting argument yields
\[
E=G=\OO\oplus\OO(-1),\quad R^2\bar f_*\OO_X=\OO(-1),
\quad0\to\bar f^*E\to\mathcal T\to\bar f^*\OO(-5)\to0. \tag{C.8}\hypertarget{tag-0099}{}\label{cdp:eq:C.8}
\]
These are the actual root-chart and cusp comparisons needed from the earlier returns; no coarse multiple fibre was replaced by its reduction in this argument.

\item Choose root sections $x$ of weight $3$, vanishing at $p_2$, and $y$ of weight $4$, vanishing at $p_1$. The Kodaira--Spencer morphism $\OO(-5)\to E\otimes G$ has, by its weights, the form
\[
 \kappa=\begin{pmatrix}0&a y\\b y&c_0x\end{pmatrix}.
\]
Its first coefficient vanishes because $H^0(\OO(5))=0$. The logarithmic period derivative at the cusp is invertible, giving $ab\ne0$. At $p_1$, the elliptic coordinate $\tau-\rho$ has a simple zero on the root chart, so its period derivative is nonzero, giving $c_0\ne0$. Rescaling the four line frames makes the three nonzero coefficients one. Cup product with this extension class gives
\[
\delta_0=(0,y,y,x)^t,\qquad
\delta_1=\begin{pmatrix}-y&0\\-x&y\end{pmatrix}:
 \OO(-5)\oplus\OO(-6)\to\OO(-1)\oplus\OO(-2).
\]
The first is a subbundle inclusion since $x,y$ have no common zero, and its cokernel is $\OO\oplus\OO(-1)\oplus\OO(2)$. The second is injective with determinant $-y^2$; its cokernel is
$Q=\OO(-1)|_{2p_1}$, a length-two module on the order-three chart. The direct-image sequence of \eqref{cdp:eq:C.8} yields
\[
 R^0\bar f_*\mathcal T=E,\quad
 R^1\bar f_*\mathcal T=\OO\oplus\OO(-1)\oplus\OO(2),\quad
 R^2\bar f_*\mathcal T=\OO(-6)\oplus Q.
\]
The last splitting follows from $H^1(Q(6))=0$; it is not designated canonical.

\item Statement \ref{cdp:st:C5a} gives the quotient $\nu_*K_{\mathrm dP_6}$. The local cusp volume form trivializes $\nu^*K_X$, and every $\nu^*M_z$ is trivial because $H^2(X,\Z)=0$ and $\Pic^0(\mathrm dP_6)=0$. Hence this quotient has cohomology only in degree two, where it is $\C$. At $p_0$, Serre duality and the conormal computation in Part A give $h^2(W,T_X|_W)=1$. Top-degree base change identifies this dimension with the residue dimension of $R^2\bar f_*T_X$. That module has generic rank one, so Nakayama's lemma and the structure theorem over $\C\{q\}$ make it free of rank one at $p_0$. Consequently its line summand is the elementary extension of $\OO(-6)$ by $\C_{p_0}$, namely $\OO(6)$, with no new cusp torsion. We have
\[
 R^0\bar f_*T_X=E,\quad
 R^1\bar f_*T_X=\OO\oplus\OO(-1)\oplus\OO(2),\quad
 R^2\bar f_*T_X=\OO(6)\oplus Q.                       \tag{C.9}\hypertarget{tag-009A}{}\label{cdp:eq:C.9}
\]

\item For a root line, removing its two stabilizer characters gives
$\rho_*\OO(n)=\OO_B(d(n))$. Averaging the finite stabilizers is exact, so there are no higher direct images of $\rho$. Its sections are equivalently monomials $x^ay^b$ with $3a+4b=n$, and Serre duality gives $h^1\OO(n)=h^0\OO(-n-7)$. At $p_1$, the invariant exponent of $Q(n)$ is $n-1$ modulo three. Its two possible coefficients have orders zero and one, so its coarse length is $\epsilon_n$ and it is killed by the coarse parameter. Apply projection formula to \eqref{cdp:eq:C.9} and then $\rho_*$. This proves \eqref{cdp:eq:C.5}, including the support and length of every torsion summand. Leray over $B$ has columns only zero and one, proving \eqref{cdp:eq:C.6}.

\item The normalization map on $R^2$ is the quotient in
\[
0\longrightarrow\OO(n-6)\longrightarrow\OO(n+6)
 \longrightarrow\C_{p_0}\longrightarrow0.
\]
Its global-section map is nonzero precisely when a section of $\OO(n+6)$ is nonzero at $p_0$. Both $x,y$ are nonzero at that point, so this occurs precisely when $n+6$ lies in the additive semigroup generated by $3,4$. That semigroup is $\{0,3,4,6,7,8,\ldots\}$, giving \eqref{cdp:eq:C.7}. The $H^1(R^1)$ part maps to zero in the cohomology of a sheaf supported at one point, by the natural Leray filtration. At $n=-1$, \eqref{cdp:eq:C.5} has only one nonzero cohomological contribution, the reduced $\C_{p_1}$ summand in degree two, while $H^3(X,\mathcal T\otimes M_{-1})=H^1(\OO(-7))=\C$. The boundary from the normalization target is an isomorphism. This proves the stated counterexample to a surjection for every twist.

\end{enumerate}
\end{proof}
